\section{Alfabet, String, Bahasa, dan Operasi Bahasa}
\label{sec:ch02-bahasa}

\subsection{Alfabet dan String}

\begin{definisi}[Alfabet, string, dan bahasa]
Alfabet \(\Sigma\) adalah himpunan berhingga yang tidak kosong dari
simbol. String adalah barisan berhingga simbol dari \(\Sigma\). String
kosong dilambangkan \(\varepsilon\) dan memiliki panjang nol. Bahasa
\(L\) atas \(\Sigma\) adalah himpunan string atas \(\Sigma\), yaitu
\(L \subseteq \Sigma^{*}\).
\end{definisi}

Himpunan seluruh string atas \(\Sigma\) dilambangkan \(\Sigma^{*}\);
himpunan string tidak kosong dilambangkan \(\Sigma^{+}\). Dengan
demikian \(\Sigma^{*} = \Sigma^{+} \cup \{\varepsilon\}\).

Panjang string \(w\) dilambangkan \(|w|\). Konkatenasi string \(u\)
dan \(v\) dilambangkan \(uv\). Pembalik string \(w\) dilambangkan
\(w^{R}\). String \(u\) adalah prefiks dari \(w\) jika terdapat \(v\)
sehingga \(w = uv\); sufiks didefinisikan secara simetris; substring
adalah potongan yang mungkin berada di tengah.

\begin{contoh}[String dan operasinya]
Misalkan \(\Sigma = \{a,b\}\) dan \(w = abba\). Maka \(|w| = 4\),
\(w^{R} = abba\), prefiks \(w\) adalah
\(\varepsilon, a, ab, abb, abba\), dan sufiksnya adalah
\(\varepsilon, a, ba, bba, abba\). String \(bb\) merupakan substring
\(w\), sedangkan \(bab\) bukan.
\end{contoh}

\subsection{Operasi atas Bahasa}

Karena bahasa adalah himpunan string, operasi himpunan berlaku:
gabungan \(L_1 \cup L_2\), irisan \(L_1 \cap L_2\), dan selisih
\(L_1 \setminus L_2\). Selain itu, terdapat operasi khas bahasa:
konkatenasi dan bintang Kleene.

\begin{definisi}[Konkatenasi dan bintang Kleene]
Konkatenasi dua bahasa didefinisikan sebagai
\[
L_1 L_2 = \{uv \mid u \in L_1,\ v \in L_2\}.
\]
Bintang Kleene \(L^{*}\) adalah himpunan konkatenasi nol atau lebih
string dari \(L\):
\[
L^{*} = \bigcup_{k \ge 0} L^{k},
\qquad L^{0} = \{\varepsilon\}.
\]
Operasi \(L^{+} = \bigcup_{k \ge 1} L^{k}\) disebut \emph{positive
closure}.
\end{definisi}

\begin{contoh}[Operasi bahasa]
Misalkan \(L_1 = \{a, ab\}\) dan \(L_2 = \{b, bb\}\). Maka
\[
L_1 L_2 = \{ab, abb, abb, abbb\} = \{ab, abb, abbb\},
\]
karena \(ab \cdot b = abb\) dan \(ab \cdot bb = abbb\), sedangkan
\(a \cdot b = ab\). Selanjutnya
\[
L_1^{*} = \{\varepsilon, a, ab, aa, aab, aba, abab, \ldots\}.
\]
Perhatikan bahwa \(L_1^{*}\) memuat string dengan jumlah \(b\) yang
sebanding dengan jumlah \(a\), tetapi tidak semua kombinasinya.
\end{contoh}

\begin{figure}[H]
\centering
\begin{tikzpicture}[node distance=6mm]
  \node[phase, text width=2.4cm] (l1) {\(L_1\)};
  \node[phase, text width=2.4cm, right=of l1] (l2) {\(L_2\)};
  \node[phase, text width=2.8cm, right=of l2] (op) {Konkatenasi};
  \node[phase, text width=2.8cm, right=of op] (hasil) {\(L_1L_2\)};
  \draw[arrow] (l1) -- (op);
  \draw[arrow] (l2) -- (op);
  \draw[arrow] (op) -- (hasil);
\end{tikzpicture}
\caption{Konkatenasi dua bahasa menghasilkan bahasa baru}
\label{fig:ch02-konkatenasi}
\end{figure}

Operasi bahasa memenuhi beberapa identitas yang berguna. Salah satunya
adalah \((L_1 \cup L_2)L_3 = L_1L_3 \cup L_2L_3\). Identitas lain
adalah \((L^{*})^{*} = L^{*}\) dan \(\varnothing^{*} =
\{\varepsilon\}\). Identitas-identitas ini akan dipakai saat
mengubah ekspresi reguler menjadi otomata dan saat membuktikan
ekuivalensi dua definisi bahasa.

\subsection{Bahasa dan Masalah Keputusan}

Setiap bahasa \(L\) dapat dipandang sebagai masalah keputusan:
diberikan string \(w\), tentukan apakah \(w \in L\). Sebuah
\emph{lexer} menjawab pertanyaan ini untuk bahasa token tertentu;
sebuah \emph{parser} menjawabnya untuk bahasa program yang lebih
kompleks. Klasifikasi bahasa menurut daya ekspresifnya, yang dibahas
pada Bagian 2.4, menentukan jenis mesin yang diperlukan untuk
menjawab pertanyaan tersebut.
